%%%%%%%%%%%%%%%%%%%%%%%%%%%%%%%%%%%%%%%%%%%%%%%%%%%%%%%%%%%%%%%%%%%%%%%%%
\makeabstract
{ % #1: Title
Swendsen-Wang Algorithm And The Invaded Cluster Algorithm
}
{ % #2: Authorship
Jeremy Lee\textsuperscript{1},
William A. Loinaz\textsuperscript{1}
}
{ % #3: Affiliations (use \\ to start a newline)
\textsuperscript{1} Department of Physics \& Astronomy, Amherst College, Amherst, MA
}
{ % #4: Purpose
 A fridge magnet suddenly losing its magnetism after reaching a certain temperature is an example of a phase transition, where a system goes from one state to another state instantly. Phase transitions are very important for understanding some of the phenomena in the world. To better understand them, we can use the $2$D Ising Model to model certain systems such as ferromagnets with phase transitions and analyze its behavior. One of the challenges is finding an appropriate algorithm to simulate the random changes in the system. In this study, we will look at cluster algorithms to see its effectiveness and efficiency in modeling phase transitions with the $2$D Ising Model. 
}
{ % #5: Methods
For this study, we will be modeling the magnetization of a flat ferromagnet with the $2$D Ising Model. The ferromagnet will be modeled with a 2D lattice of spin states, where each spin state is either up $(+1)$ or down $(-1)$. To simulate the randomness of the system, we use a cluster algorithm called the Swendsen-Wang Algorithm, where clusters are formed and flipped with a probability dependent on the temperature of the system. We measured various observables such as magnetization and compared this data with the Metropolis Algorithm data. We also used the Invaded Cluster Algorithm to calculate the critical temperature at which the phase transition occurs.
}
{ % #6: Results
We found that the Swendsen-Wang Algorithm does a great job at modeling the phase transition of the ferromagnet. In addition, we examined the equilibration and autocorrelation times of the Swendsen-Wang Algorithm, and we found that these times were significantly less at the critical temperature for the Swendsen-Wang Algorithm than the times for the Metropolis Algorithm. This demonstrates that the Swendsen-Wang Algorithm is not as affected by the effects of critical-slowing down compared to the Metropolis Algorithm. Also, our calculated value for the critical temperature by the Invaded Cluster Algorithm is $2.8736\pm0.0668$ , which is far off from the theoretical value of $2.2692$. This is likely due to the smaller lattice sizes used to carry out the measurements (which was done due to long runtime issues).
}
{ % #7: Conclusion
We found that the Swendsen-Wang Algorithm effectively models the $2$D Ising Model and significantly reduces the runtime of the simulations near the critical temperature. Also, the measurement of the critical temperature via the Invaded Cluster Algorithm can be significantly improved by optimizing the runtime of the algorithm for higher lattice sizes. 
}

%%%%%%%%%%%%%%%%%%%%%%%%%%%%%%%%%%%%%%%%%%%%%%%%%%%%%%%%%%%%%%%%%%%%%%%%%
 \newpage